\section{Grouping}
\label{sec:model}
This part of the model admits several possible encodings. A natural but
expensive formulation is a \emph{fully enumerated} encoding, where each group
explicitly stores the identifiers of all boxes it contains. Since groups may
have different lengths, this representation usually requires a two-dimensional
array with padding values for shorter groups.
\vspace{0.5cm}
\begin{lstlisting}
array[1..n_groups, 1..max_pkble_packs] of var 0..n_packs: groups;
\end{lstlisting}
\vspace{0.5cm}

Here, the value \(0\) is used as a dummy value to represent an empty slot.
This representation makes group membership explicit and easy to inspect.
However, it introduces a large number of decision variables and requires
additional constraints to enforce the structure of a valid solution. In
particular, the model must ensure that each box appears exactly once, that
non-empty positions inside each group are ordered, and that the groups are
consistent with the geometric ordering of the layer.

%\vspace{0.5cm}
%\begin{lstlisting}
%% Each pack identifier 1..n_packs appears exactly once.
%constraint
%  forall(v in 1..n_packs) (
%    count(array1d(groups), v) = 1
%  );
%
%% Non-empty entries are placed before padding values.
%constraint
%  forall(g in 1..n_groups, j in 1..max_pkble_packs-1) (
%    groups[g,j] = 0 -> groups[g,j+1] = 0
%  );
%
%% Pack identifiers are strictly increasing inside each group.
%constraint
%  forall(g in 1..n_groups, j in 1..max_pkble_packs-1) (
%    groups[g,j] != 0 /\ groups[g,j+1] != 0 ->
%      groups[g,j] < groups[g,j+1]
%  );
%\end{lstlisting}
%\vspace{0.5cm}

A second possibility is a \emph{binary assignment} encoding, where a Boolean
variable states whether pack \(i\) belongs to group \(g\). This representation
can be written as follows:

\vspace{0.5cm}
\begin{lstlisting}
array[1..n_groups, 1..n_packs] of var bool: belongs_to_group;
\end{lstlisting}
\vspace{0.5cm}

In this case, group membership is represented by the truth value of
\texttt{belongs\_to\_group[g,i]}. Although this encoding is expressive, it is
even larger than the fully enumerated representation. It also does not encode
contiguity or same-row membership directly. These properties must be
reconstructed through additional channeling constraints.

For this reason, the adopted model represents each group using its
\emph{start index} and \emph{length}. An equivalent formulation could use the
start and end indices instead. This interval-style representation is more
compact because each group is described by only two main variables. Contiguity
is implicit: once the start index and the length are known, the boxes in the
group are uniquely determined. The same-row constraint can also be expressed
directly by requiring the interval not to cross a row boundary. As a result,
the model avoids explicit membership variables and reduces the number of
constraints needed to describe valid groups.

\vspace{0.5cm}
\begin{lstlisting}
array[1..n_groups] of var 1..n_packs: group_start;
array[1..n_groups] of var 1..max_pkble_packs: group_len;
\end{lstlisting}
\vspace{0.5cm}

The adopted encoding uses two arrays: \texttt{group\_start}, which stores the
first box of each group, and \texttt{group\_len}, which stores the number of
boxes in that group.

\vspace{0.5cm}
\begin{lstlisting}
constraint group_start[1] = 1;

constraint forall(g in 1..n_groups - 1)(
  group_start[g + 1] = group_start[g] + group_len[g]
);

constraint group_start[n_groups] + group_len[n_groups] - 1 = n_packs;

constraint forall(g in 1..n_groups)(
  (group_start[g] - 1) div packs_along_x =
  (group_start[g] + group_len[g] - 2) div packs_along_x
);
\end{lstlisting}
\vspace{0.5cm}

The first three constraints define the groups as consecutive intervals over
the row-major ordering of the packs. The first group starts from pack \(1\),
and each following group starts immediately after the last pack of the previous
group. The final constraint in this block forces the last group to end at pack
\(n\). Together, these conditions ensure that all packs are covered exactly
once, without gaps or overlaps.

The fourth constraint enforces the geometric feasibility of each group. Since
the gripper can only pick boxes aligned along the \(x\)-axis, a group cannot
cross from one row to the next. The constraint therefore compares the row index
of the first pack in the group with the row index of the last pack. If the two
indices are equal, the whole interval lies on the same row.